\section{Registered Ten-Cycles}

The quotient-pentagon analysis produces twenty-four native
ten-cycles. This section proves that these are exactly the simple
ten-cycles having the registered-closure projection pattern and
determines their automorphism orbits and stabilizers.

\subsection{Definition}

Let
\[
C=(x_0,x_1,\ldots,x_9,x_0)
\]
be a simple ten-cycle in \(G_{60}\), and let
\[
\pi:G_{60}\longrightarrow G_{60}/V_4
\]
be the native Klein quotient.

We call \(C\) a \emph{registered ten-cycle} when its quotient
projection has minimal period five. Equivalently, the projected walk
traverses one quotient pentagon twice:
\[
\pi(x_{i+5})=\pi(x_i)
\]
for all \(i\), while
\[
x_{i+5}\neq x_i.
\]

The first five steps close in quotient position but not in native
state. The second traversal closes the native path.

This definition is intrinsic to the quotient map and does not depend
on a selected voltage gauge.

\subsection{Complete ten-cycle census}

The exact native graph contains
\[
1824
\]
simple ten-cycles.

Among them, exactly
\[
24
\]
project to a quotient pentagon traversed twice.

The remaining
\[
1800
\]
ten-cycles do not have this projection pattern. Their quotient
projections divide as follows:
\[
1560
\]
visit ten distinct quotient chambers, while
\[
240
\]
visit nine distinct quotient chambers.

\begin{theorem}
Among all \(1824\) simple ten-cycles in \(G_{60}\), exactly twenty-four
are registered ten-cycles.
\end{theorem}

\begin{proof}
Every one of the twelve quotient pentagons has nontrivial holonomy and
therefore lifts to two simple native ten-cycles. This gives
\[
12\cdot2=24
\]
registered cycles.

The complete simple ten-cycle census contains no additional cycle
whose quotient projection has minimal period five. Hence the
pentagon lifts are exactly the registered ten-cycles.
\end{proof}

There are therefore no false positives and no false negatives in the
projection criterion.

\subsection{Pairing above quotient pentagons}

Each quotient pentagon has a twenty-vertex preimage consisting of two
disjoint registered ten-cycles. We call these two cycles an
\emph{answering pair}.

Thus the twenty-four registered cycles are organized into
\[
12
\]
unordered pairs.

For a quotient pentagon \(Q\), write
\[
\mathcal P_Q=\{C_Q^0,C_Q^1\}
\]
for its answering pair.

The pair is intrinsic even though naming its two members requires a
choice. An automorphism may preserve the two cycles separately or
exchange them.

\begin{proposition}
The twenty-four registered ten-cycles form twelve disjoint answering
pairs, one pair above each quotient pentagon.
\end{proposition}

\subsection{One full automorphism orbit}

The full automorphism group acts transitively on the registered
ten-cycles.

\begin{theorem}
The set of twenty-four registered ten-cycles is one
\(\operatorname{Aut}(G_{60})\)-orbit.
\end{theorem}

By orbit--stabilizer,
\[
|\operatorname{Stab}(C)|
=
\frac{|\operatorname{Aut}(G_{60})|}{24}
=
\frac{480}{24}
=
20
\]
for every registered cycle \(C\).

Similarly, the twelve answering pairs form one orbit. Hence
\[
|\operatorname{Stab}(\mathcal P_Q)|
=
\frac{480}{12}
=
40.
\]

\subsection{The stabilizer of one registered cycle}

Let \(C\) be a registered ten-cycle. The setwise stabilizer acts
faithfully on the ten vertices of \(C\). Its induced action contains
the ten rotations and ten reflections of the cycle.

\begin{theorem}
The setwise stabilizer of a registered ten-cycle is the full dihedral
group
\[
\operatorname{Stab}(C)\cong D_{10},
\]
where \(D_{10}\) has order twenty.
\end{theorem}

The stabilizer contains:
\[
10
\]
orientation-preserving rotations and
\[
10
\]
orientation-reversing reflections.

Its pointwise stabilizer is trivial. Thus no nonidentity graph
automorphism fixes all ten vertices of a registered cycle.

The registered cycle therefore carries exactly the intrinsic symmetry
of an unmarked decagon.

\subsection{The stabilizer of an answering pair}

Let
\[
\mathcal P_Q=\{C_Q^0,C_Q^1\}
\]
be one answering pair.

The pair stabilizer has order forty. Exactly twenty elements preserve
the two member cycles separately, while twenty exchange them.

The exchange is generated by an involution commuting with the
dihedral action on either cycle.

\begin{theorem}
The setwise stabilizer of an answering pair is
\[
\operatorname{Stab}(\mathcal P_Q)
\cong
D_{10}\times C_2.
\]
\end{theorem}

The additional \(C_2\)-factor records whether the two lifted cycles
above the same quotient pentagon are preserved or exchanged.

The center of the pair stabilizer has order four. It contains the
central half-turn of \(D_{10}\), the answering-cycle exchange, and
their product.

\subsection{Relation to holonomy}

The two cycles in an answering pair lie above one quotient pentagon
and share its holonomy label.

Thus a \(b\)-pentagon contributes two registered cycles in the
\(b\)-family, and an \(ab\)-pentagon contributes two registered cycles
in the \(ab\)-family.

The twenty-four registered cycles therefore divide as
\[
12_b+12_{ab}.
\]

The full automorphism group exchanges these two families. Consequently
the complete set remains one orbit even though the selected native
Klein coordinates distinguish two holonomy labels.

\subsection{Registered closure as a graph property}

The registered condition can be expressed without naming a voltage
assignment:

\begin{enumerate}
    \item \(C\) is a simple ten-cycle in \(G_{60}\);
    \item its quotient projection visits exactly five chambers;
    \item the projected walk has period five;
    \item the first five native steps do not close;
    \item the second five steps complete the native cycle.
\end{enumerate}

All five conditions are preserved by graph automorphisms because the
native Klein quotient is normal and therefore invariant.

\begin{corollary}
Registered ten-cycles form an intrinsic automorphism-invariant class
of native cycles.
\end{corollary}

\subsection{Structural summary}

The registered-cycle geometry is summarized by
\[
12
\text{ quotient pentagons}
\]
\[
\longrightarrow
12
\text{ answering pairs}
\]
\[
\longrightarrow
24
\text{ registered ten-cycles}.
\]

For one cycle,
\[
\operatorname{Stab}(C)\cong D_{10}.
\]

For one answering pair,
\[
\operatorname{Stab}(\mathcal P_Q)
\cong D_{10}\times C_2.
\]

Thus a registered return carries the symmetry of a decagon, while its
answering pair adds one independent exchange bit.
